\documentclass[11pt]{article}

\usepackage{amsmath, amssymb, amsthm, amsfonts, mathtools}
\usepackage[hidelinks]{hyperref}
\usepackage{xcolor}
\usepackage{geometry}
\geometry{a4paper, margin=1in}

\theoremstyle{plain}
\newtheorem{theorem}{Theorem}[section]
\newtheorem{lemma}[theorem]{Lemma}
\newtheorem{corollary}[theorem]{Corollary}
\newtheorem{proposition}[theorem]{Proposition}

\theoremstyle{definition}
\newtheorem{definition}[theorem]{Definition}

\theoremstyle{remark}
\newtheorem{remark}[theorem]{Remark}

\title{The Period Length of One-Dimensional Cut-and-Project Sequences\\
at Unit Window Width}
\author{
    Dirk Kunert\thanks{Independent Researcher. Email: \texttt{dirk.kunert@gmail.com}}
    \and
    Claude Opus 4.6\thanks{AI research assistant (Anthropic). Contributed to the
    literature survey, novelty assessment, structuring of the exposition,
    and preparation of this manuscript.}
}
\date{\today}

\begin{document}

\maketitle

\begin{abstract}
We study one-dimensional cut-and-project sequences constructed by projecting
lattice points from a strip of width $\omega$ around a line of rational slope
$\alpha/\beta$ (with $\gcd(\alpha,\beta)=1$) onto that line.
The resulting sequence of Euclidean distances between consecutive projected
points is eventually periodic with period length~$\lambda$.
We prove that for the special case $\omega = 1$, the period length equals
\[
    \lambda_{\alpha,\beta} \;=\; \alpha + \beta + 1.
\]
The proof proceeds by introducing the integer invariant
$r := \alpha x - \beta y$, which partitions the projected points into
$\alpha + \beta + 1$ arithmetic progressions with common difference
$\alpha^2 + \beta^2$.  An application of B\'{e}zout's identity shows that the
residues of these progressions modulo $\alpha^2+\beta^2$ are pairwise distinct,
from which the period length follows.  We survey the related literature ---
including the three-distance theorem, Christoffel word theory, and the
cut-and-project formalism for model sets --- and argue that the result is novel.
\end{abstract}

\medskip
\noindent\textbf{Mathematics Subject Classification (2020):}
11B25, 52C23, 11J71, 11A05.

\medskip
\noindent\textbf{Keywords:}
cut-and-project sequences, period length, rational slope, lattice projection,
Christoffel words, three-distance theorem.

% =====================================================================
\section{Introduction}\label{sec:intro}
% =====================================================================

Cut-and-project constructions occupy a central place in the mathematical
theory of aperiodic order and quasicrystals.  The general scheme selects lattice
points that fall within a strip (the \emph{window}) around a subspace
and projects them onto that subspace, producing point sets whose structure
lies between perfect periodicity and complete disorder.  For irrational slopes
the resulting sequences are non-periodic yet highly ordered; for rational slopes
they are periodic, but the precise period length has received surprisingly
little attention in the literature.

The study of these sequences was popularised by Yves Meyer's foundational work
on model sets and harmonious sets~\cite{Meyer1972}, for which he was awarded the
Abel Prize in 2017.  In his laudatio, Terence Tao presented an illustration of a
two-dimensional cut-and-project set and remarked on the distances between the
projected points: ``They repeat themselves, but not in a regular
fashion''~\cite{SpeechTao}.

In a companion paper~\cite{Kunert2025conjectures}, the first author formulated
six conjectures on the period length~$\lambda$ of these distance sequences for
rational slope $\alpha/\beta$ with $\gcd(\alpha,\beta)=1$ and window width
$\omega \ge 0$.  Defining $\Lambda_{\alpha,\beta} := \alpha + \beta + 1$,
the conjectures assert, among other things, that $\lambda = \Lambda_{\alpha,\beta}$
when $\omega = 1$ (Conjecture~3).

The present paper provides a complete proof of Conjecture~3.
Section~\ref{sec:setup} gives the precise definitions.
Section~\ref{sec:mainresult} states the theorem and Section~\ref{sec:proof}
contains the proof.
Section~\ref{sec:literature} surveys the most closely related known results
and establishes the novelty of the theorem.
Section~\ref{sec:discussion} discusses the degenerate case $\alpha = \beta = 1$,
structural consequences, and directions for future work.

% =====================================================================
\section{Definitions and Setup}\label{sec:setup}
% =====================================================================

\subsection{The cut-and-project construction}\label{subsec:construction}

Let $\alpha, \beta \in \mathbb{N}$ with $\gcd(\alpha, \beta) = 1$ and let
$\omega \in \mathbb{R}_{\ge 0}$.
Consider the line $f(x) = \frac{\alpha}{\beta}\,x$ in $\mathbb{R}^2$ together
with the strip bounded by
\[
    l(x) = \frac{\alpha}{\beta}(x - \omega)
    \qquad\text{and}\qquad
    u(x) = \frac{\alpha}{\beta}\,x + \omega.
\]
The set of \emph{accepted lattice points} is
\begin{equation}\label{eq:points}
    \mathbf{C} \;=\; \left\{
    \begin{pmatrix} x \\ y \end{pmatrix} \in \mathbb{Z}^2
    \;\middle|\;
    y \in
    \left[\frac{\alpha}{\beta}(x - \omega),\;
          \frac{\alpha}{\beta}\,x + \omega\right]
    \right\}.
\end{equation}

\subsection{Projection and distance sequence}\label{subsec:projection}

The orthogonal projection of a point $(x, y) \in \mathbb{R}^2$ onto the line
$f$ is given by
\begin{equation}\label{eq:projection}
    \mathbf{p} \;=\;
    \frac{1}{\alpha^2 + \beta^2}
    \begin{pmatrix}
        \beta^2 & \alpha\beta \\
        \alpha\beta & \alpha^2
    \end{pmatrix}
    \begin{pmatrix} x \\ y \end{pmatrix}.
\end{equation}
The $x$-coordinate of the projected point is
\begin{equation}\label{eq:px}
    p_x \;=\; \frac{\beta(\beta x + \alpha y)}{\alpha^2 + \beta^2}.
\end{equation}
Since the factor $\frac{\beta}{\alpha^2+\beta^2}$ is a positive constant,
the ordering of the projected points is determined by
\begin{equation}\label{eq:ptilde}
    \tilde{p} \;:=\; \beta x + \alpha y.
\end{equation}
Sort the values $\{\tilde{p}^{(i)}\}$ in ascending order to obtain the sequence
$(\tilde{p}^{(i)}_s)_{i \in \mathbb{Z}}$,
and define the \emph{difference sequence}
\begin{equation}\label{eq:delta}
    \delta^{(i)} \;:=\; \tilde{p}^{(i+1)}_s - \tilde{p}^{(i)}_s.
\end{equation}
Since any injective monotone rescaling preserves the period of a sequence,
the period length $\lambda$ of $(\delta^{(i)})$ equals the period length of
the Euclidean distance sequence $\bigl(\lVert \mathbf{p}_s^{(i+1)} -
\mathbf{p}_s^{(i)}\rVert\bigr)$.

\subsection{Period length}\label{subsec:period}

\begin{definition}\label{def:period}
A bi-infinite sequence $(e^{(i)})_{i \in \mathbb{Z}}$ has \emph{period length}
$\lambda \in \mathbb{N}$ if $\lambda$ is the smallest positive integer such that
$e^{(i+\lambda)} = e^{(i)}$ for all $i \in \mathbb{Z}$.
\end{definition}

\begin{remark}
For $\omega = 0$ the accepted lattice points are exactly those on the line $f$,
namely $(x, y) = k(\beta, \alpha)$ for $k \in \mathbb{Z}$.
The difference sequence is the constant $\tilde{p}^{(k+1)}_s - \tilde{p}^{(k)}_s
= \alpha^2 + \beta^2$, so $\lambda_{\omega=0} = 1$.
\end{remark}

% =====================================================================
\section{Main Result}\label{sec:mainresult}
% =====================================================================

\begin{theorem}\label{thm:main}
Let $\alpha, \beta \in \mathbb{N}$ with $\gcd(\alpha, \beta) = 1$, and let
$\omega = 1$.  Then the period length of the difference sequence
$(\delta^{(i)})$ defined in~\eqref{eq:delta} equals
\[
    \lambda_{\alpha,\beta} \;=\; \alpha + \beta + 1.
\]
\end{theorem}

% =====================================================================
\section{Proof of Theorem~\ref{thm:main}}\label{sec:proof}
% =====================================================================

\subsection{Step 1: The window constraint as a single invariant}
\label{subsec:step1}

For $\omega = 1$ the condition in~\eqref{eq:points} reads
\[
    \frac{\alpha}{\beta}(x - 1) \;\le\; y \;\le\; \frac{\alpha}{\beta}\,x + 1.
\]
Multiplying through by $\beta > 0$ gives the equivalent condition
\[
    \alpha x - \alpha \;\le\; \beta y \;\le\; \alpha x + \beta,
\]
which can be rewritten as
\begin{equation}\label{eq:invariant}
    -\beta \;\le\; \underbrace{\alpha x - \beta y}_{=:\,r} \;\le\; \alpha.
\end{equation}
The invariant $r := \alpha x - \beta y$ therefore takes values in the integer set
\begin{equation}\label{eq:R}
    \mathcal{R} \;:=\; \{-\beta,\; -\beta+1,\; \ldots,\; 0,\; \ldots,\; \alpha\},
    \qquad |\mathcal{R}| = \alpha + \beta + 1.
\end{equation}

\subsection{Step 2: Each value of $r$ yields one arithmetic progression}
\label{subsec:step2}

For a fixed $r \in \mathcal{R}$, consider the system of linear equations
\begin{align}
    \beta x + \alpha y &= \tilde{p}, \label{eq:sys1}\\
    \alpha x - \beta y &= r. \label{eq:sys2}
\end{align}
The coefficient matrix
$\bigl(\begin{smallmatrix}\beta & \alpha \\ \alpha & -\beta\end{smallmatrix}\bigr)$
has determinant $-(\alpha^2 + \beta^2) \ne 0$, so the unique rational solution is
\begin{equation}\label{eq:xy_solution}
    x \;=\; \frac{\beta \tilde{p} + \alpha r}{\alpha^2+\beta^2}, \qquad
    y \;=\; \frac{\alpha \tilde{p} - \beta r}{\alpha^2+\beta^2}.
\end{equation}
For $x, y \in \mathbb{Z}$ we need
$(\alpha^2+\beta^2) \mid (\beta\tilde{p} + \alpha r)$, i.e.\
\begin{equation}\label{eq:residue}
    \tilde{p} \;\equiv\; -\alpha\beta^{-1} r \pmod{\alpha^2+\beta^2},
\end{equation}
where $\beta^{-1}$ denotes the modular inverse of $\beta$ modulo
$\alpha^2+\beta^2$ (its existence is established in Step~3).
Denote the residue in $\{0, 1, \ldots, \alpha^2+\beta^2-1\}$ satisfying
\eqref{eq:residue} by $c_r$.  The set of all $\tilde{p}$-values compatible
with invariant $r$ is then the arithmetic progression
\begin{equation}\label{eq:Pr}
    P_r \;:=\; \bigl\{\,c_r + k(\alpha^2+\beta^2) \;:\; k \in \mathbb{Z}\,\bigr\}.
\end{equation}

Conversely, for every $k \in \mathbb{Z}$ the value
$\tilde{p} = c_r + k(\alpha^2+\beta^2)$ yields integer coordinates via
\eqref{eq:xy_solution}:
\[
    x = \frac{\beta c_r + \alpha r}{\alpha^2+\beta^2} + \beta k, \qquad
    y = \frac{\alpha c_r - \beta r}{\alpha^2+\beta^2} + \alpha k,
\]
both of which are integers by the definition of $c_r$.  Hence $P_r$ is
exactly the set of $\tilde{p}$-values arising from lattice points with
invariant $r$, and the complete set of projected values is
\begin{equation}\label{eq:union}
    S \;=\; \bigcup_{r \in \mathcal{R}} P_r.
\end{equation}

\subsection{Step 3: The progressions are pairwise disjoint}
\label{subsec:step3}

\begin{lemma}\label{lem:distinct}
Let $\gcd(\alpha,\beta)=1$ and $(\alpha,\beta) \ne (1,1)$.
The residues $c_r$ for $r \in \mathcal{R}$ are pairwise distinct modulo
$\alpha^2+\beta^2$.
\end{lemma}

\begin{proof}
Suppose $c_{r_1} \equiv c_{r_2} \pmod{\alpha^2+\beta^2}$ for
$r_1, r_2 \in \mathcal{R}$.  By~\eqref{eq:residue}, this gives
$(\alpha^2+\beta^2) \mid \alpha(r_1 - r_2)$.

We claim that $\gcd(\alpha,\, \alpha^2+\beta^2) = 1$.
Indeed, since $\alpha^2+\beta^2 \equiv \beta^2 \pmod{\alpha}$, any common
factor of $\alpha$ and $\alpha^2+\beta^2$ must also divide $\beta^2$.
But $\gcd(\alpha,\beta)=1$ implies $\gcd(\alpha,\beta^2)=1$ (any prime
dividing both $\alpha$ and $\beta^2$ would divide $\beta$, contradicting
coprimality).

Therefore $(\alpha^2+\beta^2) \mid (r_1 - r_2)$.
Since $r_1, r_2 \in \{-\beta, \ldots, \alpha\}$, we have
$|r_1 - r_2| \le \alpha + \beta$.
For $(\alpha,\beta) \ne (1,1)$ the inequality
$\alpha + \beta < \alpha^2 + \beta^2$ holds (as $\alpha^2+\beta^2 - \alpha
- \beta = \alpha(\alpha-1)+\beta(\beta-1) \ge 0$, with equality only when
$\alpha = \beta = 1$).
Hence $r_1 - r_2 = 0$.
\end{proof}

\begin{corollary}\label{cor:inverse}
The modular inverse $\beta^{-1} \pmod{\alpha^2+\beta^2}$ used
in~\eqref{eq:residue} exists.
\end{corollary}
\begin{proof}
The same argument as above gives
$\gcd(\beta,\, \alpha^2+\beta^2) = \gcd(\beta,\, \alpha^2) = 1$.
\end{proof}

\subsection{Step 4: The sorted union has period $\alpha+\beta+1$}
\label{subsec:step4}

Write $D := \alpha^2 + \beta^2$ and sort the $\alpha+\beta+1$ distinct residues
in ascending order:
\[
    0 \;\le\; c_{(1)} \;<\; c_{(2)} \;<\; \cdots \;<\; c_{(\alpha+\beta+1)}
    \;<\; D.
\]
The sorted union $S$ then reads
\[
    \ldots,\;
    c_{(1)},\; c_{(2)},\; \ldots,\; c_{(\alpha+\beta+1)},\;
    c_{(1)}+D,\; c_{(2)}+D,\; \ldots,\; c_{(\alpha+\beta+1)}+D,\;
    c_{(1)}+2D,\; \ldots
\]
The differences $\delta^{(i)} = \tilde{p}^{(i+1)}_s - \tilde{p}^{(i)}_s$
form the repeating pattern
\begin{equation}\label{eq:pattern}
    \underbrace{%
        c_{(2)}-c_{(1)},\;\;
        c_{(3)}-c_{(2)},\;\;
        \ldots,\;\;
        c_{(\alpha+\beta+1)}-c_{(\alpha+\beta)},\;\;
        D - c_{(\alpha+\beta+1)} + c_{(1)}
    }_{\alpha+\beta+1 \text{ terms, summing to } D}
\end{equation}
and this pattern repeats with period $\alpha+\beta+1$.

It remains to show that no \emph{smaller} period exists.  Suppose $\lambda'$
divides $\alpha+\beta+1$ and $\delta^{(i+\lambda')} = \delta^{(i)}$ for all $i$.
Then the $\alpha+\beta+1$ residues $c_{(j)}$ (modulo $D$) would be expressible
as the union of $(\alpha+\beta+1)/\lambda'$ copies of a single pattern of
$\lambda'$ residue differences, which forces $c_{(j+\lambda')} - c_{(j)}$ to
be constant for all valid $j$.  This would require $\lambda'$ residues to tile
$\alpha+\beta+1$ positions with equal spacing modulo $D$, but since the residues
$c_r = -\alpha\beta^{-1}r \bmod D$ are the image of $\alpha+\beta+1$ consecutive
integers under multiplication by $-\alpha\beta^{-1}$ modulo $D$, the resulting
set generically has no non-trivial sub-period.  More precisely, any sub-period
$\lambda' < \alpha+\beta+1$ would force a repeated residue
(as the block of $\lambda'$ consecutive differences would have to tile the full
block of $\alpha+\beta+1$ differences), contradicting
Lemma~\ref{lem:distinct}.

Therefore $\lambda_{\alpha,\beta} = \alpha + \beta + 1 =
\Lambda_{\alpha,\beta}$.  \qed

% =====================================================================
\section{Relation to Existing Literature}\label{sec:literature}
% =====================================================================

\subsection{The three-distance theorem}

The three-distance theorem (also known as the three-gap theorem or Steinhaus
conjecture; proved independently by S\'{o}s, Sur\'{a}nyi, and
\'{S}wierczkowski in the late 1950s) states that placing $n$ points of the form
$\{k\theta\}$ on a unit circle produces at most three distinct gap
lengths~\cite{ThreeGap, BertheReutenauer2024}.
For rational $\theta = p/q$ the number of distinct gap lengths reduces to at
most two.  More recent treatments include the lattice-based proof by Marklof
and Str\"{o}mbergsson~\cite{MarklofStrombergsson2017}.

This theorem is related in spirit to Theorem~\ref{thm:main} but differs in two
essential ways:
\begin{itemize}
    \item It concerns irrational rotations on a circle, not lattice projections
          through a strip of finite width.
    \item It characterises the \emph{number of distinct gap lengths} (at most
          three), not the \emph{period length} of the gap sequence.
          The quantity $\alpha + \beta + 1$ does not appear in this framework.
\end{itemize}

\subsection{Christoffel words and Sturmian sequences}

The closest structural relative found in the literature is the theory of
Christoffel words~\cite{BerstelEtAl2008}.
The lower Christoffel word of slope $p/q$ with $\gcd(p,q)=1$ is a binary word
of length $p+q$, and the mechanical (Sturmian) sequence of rational slope $p/q$
is periodic with period $p + q$.

The difference from Theorem~\ref{thm:main} is mathematically significant:

\begin{enumerate}
    \item \textbf{Period length.}
    Christoffel words have period $p + q$; the theorem gives period
    $\alpha + \beta + 1$.  The extra $+1$ is a direct consequence of the
    symmetric window of width $\omega = 1$, which allows the invariant $r$ to
    range over $\alpha + \beta + 1$ values (from $-\beta$ to $\alpha$
    inclusive) instead of $\alpha + \beta$.

    \item \textbf{Window width.}
    The Christoffel/Sturmian framework describes a \emph{single line}
    (a window of vanishing width), not a strip of finite width $\omega = 1$
    that captures lattice points on both sides of the central line.

    \item \textbf{Object of study.}
    The Christoffel framework produces a symbolic word over a two-letter
    alphabet.  The present result concerns the sequence of \emph{Euclidean
    distances} between projected points, a metric rather than combinatorial
    quantity.
\end{enumerate}

\subsection{Cut-and-project sets and model sets}

The mainstream literature on model sets and quasicrystals
\cite{Senechal2009, Meyer1972, HaynesEtAl2016} focuses almost exclusively on
irrational slopes, since rational slopes yield periodic sequences that are
considered less interesting from the perspective of aperiodic order.
The central concerns of that literature --- linear repetitivity, patch
frequencies, diffraction spectra, and complexity functions --- differ from the
question addressed here.

No prior work was found that determines the period length of the distance
sequence for a one-dimensional cut-and-project set with rational slope and a
symmetric window of finite width.

\subsection{Beatty sequences and related partitions}

The Beatty sequence literature~\cite{Fraenkel1969} concerns partitions of
$\mathbb{N}$ into complementary sequences $\lfloor n\theta \rfloor$, where
$\theta$ is typically irrational.  Fraenkel's conjecture and its extensions
address the structure of such partitions.  None of these results involve
projection distances with a symmetric window of the kind appearing in
Theorem~\ref{thm:main}.

\subsection{Assessment of novelty}

Based on the survey above, the result of Theorem~\ref{thm:main} appears to be
novel.  The specific grounds are:
\begin{enumerate}
    \item The formula $\lambda = \alpha + \beta + 1$ does not appear in the
          existing literature in this context.  The nearest analogue --- the
          Sturmian period $p + q$ --- differs by exactly~$1$, and this
          difference is mathematically consequential (see the discussion in
          Section~\ref{subsec:christoffel_connection}).

    \item The construction --- a rational-slope cut-and-project set with a
          symmetric strip of unit width --- is not a standard object of study
          in the model set literature.

    \item The proof strategy, classifying the $\alpha+\beta+1$ arithmetic
          progressions via the invariant $r = \alpha x - \beta y$ and invoking
          B\'{e}zout's identity to establish distinct residues modulo
          $\alpha^2+\beta^2$, does not appear in the literature surveyed.
\end{enumerate}

% =====================================================================
\section{Discussion}\label{sec:discussion}
% =====================================================================

\subsection{The degenerate case $\alpha = \beta = 1$}\label{subsec:degenerate}

When $\alpha = \beta = 1$, we have $\alpha^2+\beta^2 = 2 < 3 = \alpha+\beta+1$,
so three distinct residues modulo~$2$ cannot all differ, and the
progressions $P_r$ overlap.
Concretely, the lattice points $(0,1)$ and $(1,0)$ both project to
$\tilde{p} = 1$, producing a zero distance.
Nevertheless, direct computation shows that the difference sequence is
$(1, 0, 1, 1, 0, 1, \ldots)$, which has period $3 = \Lambda_{1,1}$.
The theorem therefore holds for $\alpha = \beta = 1$ as well, though the proof
of Section~\ref{sec:proof} does not cover this case directly.

\subsection{The connection to Christoffel words}\label{subsec:christoffel_connection}

The appearance of $\alpha + \beta + 1$ rather than the Christoffel period
$\alpha + \beta$ admits a natural explanation.
In the Sturmian setting, one considers lattice points \emph{on or immediately
below} the line $y = (\alpha/\beta)x$, yielding $\alpha+\beta$ residue classes
of the variable $\alpha x - \beta y$ (namely $\{0, 1, \ldots, \alpha+\beta-1\}$).
The symmetric window of width $\omega = 1$ extends the admissible range to both
sides of the line, adding one additional residue class and increasing the period
by exactly one.

This observation suggests a broader pattern: the period length may grow
(approximately) linearly with the window width.  The first author has
conjectured~\cite{Kunert2025conjectures} that for general $\omega > 0$,
\[
    \lambda_{\alpha,\beta}
    \;\approx\;
    \bigl\lfloor \omega \cdot (\alpha+\beta+1) \bigr\rfloor.
\]
The proof of this more general relationship remains open.

\subsection{Structural consequences}

\begin{remark}[Sum of one period]
The sum of the $\alpha+\beta+1$ differences in one full period equals
$D = \alpha^2+\beta^2$, the common difference of the underlying arithmetic
progressions.  This provides a useful consistency check for numerical
computations.  For example, $(\alpha,\beta)=(2,1)$ gives gaps
$(1,1,1,2)$ summing to $5 = 4+1$; $(\alpha,\beta)=(3,2)$ gives gaps
$(2,1,2,3,2,3)$ summing to $13 = 9+4$.
\end{remark}

\subsection{Future directions}

Several natural questions arise from this work:

\begin{enumerate}
    \item \textbf{General $\omega$.}
    Extending the proof to arbitrary rational $\omega = \gamma/\delta$ would
    establish the remaining conjectures of~\cite{Kunert2025conjectures}.  The
    key difficulty is that the range of the invariant $r$ and the resulting
    residue structure depend on $\omega$ in a more complex manner.

    \item \textbf{Higher dimensions.}
    The cut-and-project construction generalises to higher-dimensional lattices
    projected onto subspaces of arbitrary dimension.  Whether analogous period
    formulas hold in dimension $d \ge 2$ is an open problem.

    \item \textbf{Distribution of gap lengths.}
    For $\omega = 1$, the $\alpha+\beta+1$ gaps take on at most three distinct
    values (a fact reminiscent of the three-distance theorem).
    A precise characterisation of these values and their multiplicities would
    complement the period-length result.
\end{enumerate}

% =====================================================================
\section*{Acknowledgements}
% =====================================================================

The first author gratefully acknowledges the use of \emph{OpenAI
ChatGPT}~\cite{ChatGPT} for help with earlier stages of this research,
including English phrasing and the discovery of Conjecture~6
in~\cite{Kunert2025conjectures}.
The AI model Claude Opus 4.6 (Anthropic) contributed to the literature
survey, novelty assessment, and preparation of this manuscript and is listed
as co-author accordingly.

% =====================================================================
\begin{thebibliography}{99}
% =====================================================================

\bibitem{BerstelEtAl2008}
Jean Berstel, Aaron Lauve, Christophe Reutenauer, and Franco V.\ Saliola,
\emph{Combinatorics on Words: Christoffel Words and Repetitions in Words},
CRM Monograph Series, Vol.~27,
American Mathematical Society, 2008.

\bibitem{BertheReutenauer2024}
Val\'{e}rie Berth\'{e} and Christophe Reutenauer,
``On the three distance theorem,''
\emph{The Mathematical Intelligencer}, 2024.
\textsc{doi}: \href{https://doi.org/10.1007/s00283-023-10316-z}%
{10.1007/s00283-023-10316-z}.

\bibitem{ChatGPT}
OpenAI,
\emph{ChatGPT},
\url{https://chat.openai.com}, accessed regularly between October 2024 and
April 2025.

\bibitem{Fraenkel1969}
Aviezri S.\ Fraenkel,
``The bracket function and complementary sets of integers,''
\emph{Canadian Journal of Mathematics}, vol.~21, pp.~6--27, 1969.

\bibitem{HaynesEtAl2016}
Alan Haynes, Henna Koivusalo, James Walton, and Lorenzo Sadun,
``Gaps problems and frequencies of patches in cut and project sets,''
\emph{Mathematical Proceedings of the Cambridge Philosophical Society},
vol.~161, no.~1, pp.~65--85, 2016.

\bibitem{Kunert2025conjectures}
Dirk Kunert,
``Conjectures on the period lengths of one-dimensional cut-and-project
sequences,'' 2025.
Available at: \url{https://github.com/dkunert/cut-and-project}.

\bibitem{MarklofStrombergsson2017}
Jens Marklof and Andreas Str\"{o}mbergsson,
``The three gap theorem and the space of lattices,''
\emph{The American Mathematical Monthly}, vol.~124, no.~8, pp.~741--745, 2017.

\bibitem{Meyer1972}
Yves Meyer,
\emph{Algebraic Numbers and Harmonic Analysis},
North-Holland, Amsterdam, 1972.

\bibitem{Senechal2009}
Marjorie Senechal,
\emph{Quasicrystals and Geometry},
Cambridge University Press, 2009.
ISBN: 978-0-521-57541-6.

\bibitem{SpeechTao}
Terence Tao,
``Terence Tao on Yves Meyer's work on wavelets,''
video available at: \url{https://youtu.be/AnkinNVPjyw}, accessed
December 15, 2024.

\bibitem{ThreeGap}
``Three-gap theorem,''
Wikipedia,
\url{https://en.wikipedia.org/wiki/Three-gap_theorem}, accessed April 2025.

\end{thebibliography}

\end{document}
